% !TeX root = ../../Book4.tex
% 纲要标题：### 16.2 超同调谱序列
% 纲要标签：b4:sec:1502
% 纲要位置：计划纲要.md，第 3630 行。

\section{超同调谱序列}
\label{b4:sec:1502}

对复形逐项取内射消解，还必须使原微分的提升严格满足平方为零。Cartan--Eilenberg 消解同时照顾项、圈、边界和同调，才能支持超上同调计算。前两小节给出消解的定义、存在性与比较证明，也是\secref{b4:sec:1704}建立有界消解模型所需的部分；本节末的超上同调谱序列另以\chapref{b4:ch:14}为基础。

% 纲要内容：
%
% * 复形上的函子
% * 消解双复形
% * 由对象同调通向总同调
%

\subsection{复形上的函子}
\label{b4:subsec:1502:01}

设 \(F:\mathcal A\to\mathcal B\) 是交换范畴间的加性左正合函子，\(\mathcal A\) 有足够内射对象。对单个对象，取内射消解后作用 \(F\) 已定义 \(R^nF\)。对复形 \(C^\bullet\)，仅逐项独立取内射消解还不够：原来的微分也必须提升，并且要保持平方为零。若分别选择 \(C^a\to I^{a,\bullet}\) 及微分的比较提升 \(\widetilde{\mathrm d}^{\,a}:I^{a,\bullet}\to I^{a+1,\bullet}\)，则 \(\widetilde{\mathrm d}^{\,a+1}\widetilde{\mathrm d}^{\,a}\) 与零链映射都提升原来的零映射 \(\mathrm d_C^{a+1}\mathrm d_C^a\)。由\cref{b4:prop:comparison-homotopy-unique}，二者链同伦；但同伦于零不等于严格为零，因而这些提升未必组成双复形。

\ProseParagraphBreak
解决方法是同时消解项、余圈、余边和上同调，使每行的核与像受同一构造控制。以下只处理有下界的上链复形。先在原微分次数记为 \(a\)、消解次数记为 \(b\) 的交换方格中构造，最后按 \(\mathrm{d}_v=(-1)^a\partial_v\)、\(\mathrm{d}_h=\partial_h\) 转成反交换方格。

\begin{definition}\label{b4:def:ce-resolution}
设 \(\mathcal A\) 是有足够内射对象的交换范畴，\(C\) 是其中的有下界上链复形。给定带增广 \(C^a\to I^{a,0}\) 的交换双复形 \((I^{a,b},\partial_h,\partial_v)\)，其中 \(b\ge0\)。若对每个 \(a\)，以下四个增广纵列都是内射消解：
\[
 C^a\to I^{a,\bullet},\quad
 B^a(C)\to B_h^a(I^{\bullet,\bullet}),\quad
 Z^a(C)\to Z_h^a(I^{\bullet,\bullet}),\quad
 H^a(C)\to H_h^a(I^{\bullet,\bullet}).
\]
并且 \(C\) 在某个次数以下为零时，\(I\) 也在该横次数以下为零，则称 \(I\) 为 \(C\) 的\term[ce-xiaojie]{Cartan--Eilenberg 消解}[Cartan--Eilenberg resolution]。这里 \(B_h,Z_h,H_h\) 分别按横微分取余边、余圈和上同调。
\end{definition}

之所以要求这些纵列的各项都内射，是为了使每一横行中的
\[
 0\to B_h^a\to Z_h^a\to H_h^a\to0,
 \qquad 0\to Z_h^a\to I^{a,b}\to B_h^{a+1}\to0
\]
都分裂：第一列的左端 \(B_h^a\) 内射，第二列的左端 \(Z_h^a\) 内射，故相应包含都有收缩。这使用的是左端的内射性；仅有右端内射不能保证分裂。任意加性函子保持分裂短正合列，横向上同调与 \(F\) 才能交换。

\subsection{消解双复形}
\label{b4:subsec:1502:02}

\begin{lemma}\label{b4:lem:ce-injective-rows}
设 \(\mathcal A\) 是有足够内射对象的交换范畴。每个 \(\mathcal A\) 中的有下界上链复形 \(C\) 都能嵌入一个有相同下界的复形 \(J\)，使 \(J^a,B^a(J),Z^a(J),H^a(J)\) 均内射，且 \(C\to J\) 在项、余边及上同调上均为单态射。任意满足这些内射性条件的有下界复形 \(J\) 都有如下延拓性质：若有下界复形映射 \(X\to Y\) 在项、余边及上同调上均为单态射，则每个复形映射 \(X\to J\) 都能延拓为 \(Y\to J\)。
\end{lemma}
\begin{proof}
记 \(q_C^a:C^a\twoheadrightarrow C^a/B^a(C)\) 为商态射，\(\beta_C^a:C^a\twoheadrightarrow B^{a+1}(C)\) 为微分的像分解中的满态射。余圈包含诱导单态射
\(H^a(C)\hookrightarrow C^a/B^a(C)\)。选择单态射
\(u^a:B^a(C)\hookrightarrow U^a\)、\(v^a:H^a(C)\hookrightarrow V^a\)，其中 \(U^a,V^a\) 内射；对象为零时取零。由内射性，分别将它们延拓为
\[
\alpha^a:C^a\longrightarrow U^a,\qquad
\eta^a:C^a/B^a(C)\longrightarrow V^a.
\]
构造 \(J\) 时要分别保存这一度的余边、这一度的上同调，以及将由微分送到下一度的余边。\(U^a\) 承担 \(B^a(C)\) 的内射嵌入，\(V^a\) 承担 \(H^a(C)\) 的内射嵌入；第三份则必须使用下一度的 \(U^{a+1}\)，才能接收 \(\beta_C^a:C^a\twoheadrightarrow B^{a+1}(C)\) 的像。因此令
\[
J^a=U^a\oplus V^a\oplus U^{a+1},\qquad
\mathrm{d}_J^a=
\begin{pmatrix}
0&0&1_{U^{a+1}}\\0&0&0\\0&0&0
\end{pmatrix}.
\]
这三项的分工可写在微分旁边：
\[
\begin{aligned}
J^a&=\underbrace{U^a}_{B^a(C)\text{ 的嵌入}}
\oplus\underbrace{V^a}_{H^a(C)\text{ 的嵌入}}
\oplus\underbrace{U^{a+1}}_{B^{a+1}(C)\text{ 的嵌入}},\\
\mathrm{d}_J^a(x,y,z)&=(z,0,0)
\in U^{a+1}\oplus V^{a+1}\oplus U^{a+2}.
\end{aligned}
\]
微分只把当前的第三块送到下一度的第一块，因为它记录的正是下一度余边；当前余边与上同调代表都应被微分杀掉，所以其余块为零。再次微分时，这个像已位于第一块，因而自动为零。
矩阵表示有限双积之间的态射，其靶为 \(J^{a+1}\)。由此
\(B^a(J)=U^a\)、\(Z^a(J)=U^a\oplus V^a\)、\(H^a(J)=V^a\)，均内射。以双积的三个投影规定态射
\[
j^a=\begin{pmatrix}
\alpha^a\\ \eta^a q_C^a\\ u^{a+1}\beta_C^a
\end{pmatrix}:C^a\longrightarrow J^a.
\]
由于 \(\alpha^{a+1}\mathrm{d}_C^a=u^{a+1}\beta_C^a\)、
\(q_C^{a+1}\mathrm{d}_C^a=0\) 和
\(\beta_C^{a+1}\mathrm{d}_C^a=0\)，直接比较三个投影就有
\(j^{a+1}\mathrm{d}_C^a=\mathrm{d}_J^a j^a\)。

为证明 \(j^a\) 单，取其核 \(\kappa:K\hookrightarrow C^a\)。第三个投影给出 \(u^{a+1}\beta_C^a\kappa=0\)，由 \(u^{a+1}\) 单得到 \(\mathrm{d}_C^a\kappa=0\)，故 \(\kappa\) 经 \(Z^a(C)\) 因子化。将这一因子投影到 \(H^a(C)\)，第二个投影及 \(v^a\) 的单射性使该投影为零，因此 \(\kappa\) 又经 \(B^a(C)\) 因子化。最后第一个投影在 \(B^a(C)\) 上就是单态射 \(u^a\)，所以该因子为零，\(K=0\)。\(j\) 在余边和上同调上诱导的态射分别是 \(u^a,v^a\)，故也单。若 \(C^a=0\) 对 \(a<m\) 成立，则还有 \(B^m(C)=0\)；因而 \(U^m=0\)，连同较低次数所取的零对象，保证 \(J^{m-1}=0\)，并使 \(J\) 有相同下界。

下面证明延拓性质。设 \(i:X\hookrightarrow Y\) 满足题设。在模范畴中，若 \(x\in X^a\) 的像在 \(Y^a\) 中成为余边，则 \(i(\mathrm d_Xx)=\mathrm d_Yi(x)=0\)，项上的单射性使 \(x\) 为余圈。再由 \(H^a(i)([x])=0\) 和上同调单射性，得到 \([x]=0\)，故 \(x\) 原本就是 \(X\) 中的余边。一般交换范畴中，用拉回来表示这份交。令
\(P=X^a\times_{Y^a}B^a(Y)\)，并记其到 \(X^a\) 的投影为 \(p\)。链映射关系给出
\(i^{a+1}\mathrm{d}_X^ap=0\)；\(i^{a+1}\) 单使 \(p\) 经 \(Z^a(X)\) 因子化。该因子在 \(H^a(Y)\) 中的像为零，而 \(H^a(i)\) 单，故它在 \(H^a(X)\) 中也为零，进一步经 \(B^a(X)\) 因子化。反过来，\(i\) 把 \(B^a(X)\) 映入 \(B^a(Y)\)，所以这两个因子给出 \(P\simeq B^a(X)\) 作为 \(X^a\) 的子对象。核与余核的因子化因而给出
\[
\ker\bigl(X^a/B^a(X)\longrightarrow Y^a/B^a(Y)\bigr)
\simeq P/B^a(X)=0.
\]
这证明所需商态射是单态射。

对于任意满足引理内射性条件的 \(J\)，两条短正合列
\[
\begin{gathered}
0\longrightarrow B^a(J)\longrightarrow Z^a(J)
\longrightarrow H^a(J)\longrightarrow0,\\
0\longrightarrow Z^a(J)\longrightarrow J^a
\longrightarrow B^{a+1}(J)\longrightarrow0.
\end{gathered}
\]
因左端内射而分裂。选择分裂就得到
\[
J^a\simeq B^a(J)\oplus H^a(J)\oplus B^{a+1}(J),
\]
且在这些坐标下，微分仍是刚才的标准矩阵：它投影到 \(B^{a+1}(J)\)，再通过下一次的余边包含映入 \(J^{a+1}\)。因此只需对这种标准式证明延拓。

记 \(U^a=B^a(J)\)、\(V^a=H^a(J)\)。复形态射 \(f:X\to J\) 的中间投影杀掉 \(B^a(X)\)，故因子化为 \(X^a/B^a(X)\to V^a\)；其第三个投影由链映射关系等于 \(f_U^{a+1}\mathrm{d}_X^a\)，其中 \(f_U^a:X^a\to U^a\) 是第一个投影。所求延拓也没有第三个分量的自由选择：由 \(\mathrm d_J^a(x,y,z)=(z,0,0)\)，链映射条件强制 \(\pi_3g^a=\pi_1g^{a+1}\mathrm d_Y^a\)。因此利用 \(U^a,V^a\) 的内射性，沿 \(X^a\hookrightarrow Y^a\) 及刚证明的商单态射延拓前两个投影，分别记为 \(g_U^a\) 和 \(g_V^a\)，再定义
\[
g^a=\begin{pmatrix}
g_U^a\\g_V^a q_Y^a\\g_U^{a+1}\mathrm{d}_Y^a
\end{pmatrix}:Y^a\longrightarrow J^a.
\]
第三投影的规定与中间投影杀掉余边的性质保证
\(g^{a+1}\mathrm{d}_Y^a=\mathrm{d}_J^ag^a\)。三个投影在 \(X^a\) 上都恢复 \(f^a\)，故 \(g\) 就是所需延拓。
\end{proof}

\begin{theorem}\label{b4:thm:ce-resolution}
设 \(\mathcal A\) 是有足够内射对象的交换范畴。每个 \(\mathcal A\) 中的有下界上链复形都有 Cartan--Eilenberg 消解。对任意两个这样的复形 \(C,D\) 及其消解，复形映射 \(C\to D\) 可提升为消解之间的交换双复形映射；两种提升诱导的总复形映射，在作用任意到交换范畴的加性函子后仍链同伦。因此总化后的上同调，以及\cref{b4:thm:hypercohomology-spectral-sequences}中两套谱序列从第二页起的自然识别，都不依赖消解选择。
\end{theorem}
\begin{proof}
先考察复形短正合列 \(0\to X\xrightarrow{i}Y\to K\to0\)，假设 \(H^a(i)\) 对每个 \(a\) 都单。由\cref{b4:thm:cohomology-long-exact}，所有连接态射
\(\delta^a:H^a(K)\to H^{a+1}(X)\) 为零，因而有上同调短正合列。将次数 \(a,a+1\) 的项短正合列及其微分组成交换图，蛇引理给出余圈列的连接态射
\(Z^a(K)\to X^{a+1}/B^{a+1}(X)\)。复形关系 \(\mathrm d^2=0\) 使其像落在 \(H^{a+1}(X)\) 这一子对象中；它在 \(B^a(K)\) 上为零，所以下降后正是上同调长正合列的 \(\delta^a\)。于是该态射因子化为
\[
Z^a(K)\twoheadrightarrow H^a(K)
\xrightarrow{\delta^a}H^{a+1}(X)
\hookrightarrow X^{a+1}/B^{a+1}(X),
\]
故为零。这证明 \(Z^a(Y)\to Z^a(K)\) 满，并得到
\[
0\longrightarrow Z^a(X)\longrightarrow Z^a(Y)\longrightarrow Z^a(K)\longrightarrow0.
\]
在模范畴中，这正是说 \(K\) 中的余圈可以提升为 \(Y\) 中的余圈：任取提升的微分落在 \(X\) 中，其上同调类是连接态射的值；该值为零时，减去 \(X\) 中一个适当元素便消去微分。

再对 \(X,Y\) 的两条短正合列
\(0\to Z^a(-)\to(-)^a\to B^{a+1}(-)\to0\)
使用蛇引理。项上的单射使余边映射也单，而刚证明的余圈正合性将两个余核分别识别为 \(Z^a(K)\) 与 \(K^a\)。蛇引理因此给出
\[
0\longrightarrow Z^a(K)\longrightarrow K^a
\longrightarrow\operatorname{coker}\bigl(B^{a+1}(X)\to B^{a+1}(Y)\bigr)
\longrightarrow0.
\]
与 \(B^{a+1}(K)=\operatorname{coker}(Z^a(K)\hookrightarrow K^a)\) 比较，便得到余边短正合列。至此，取余复形同时保持了项、余圈、余边和上同调四类短正合列。

令 \(K_0=C\)。逐次用上一引理选择 \(K_b\hookrightarrow I^{\bullet,b}\)，并令
\(K_{b+1}=\operatorname{coker}(K_b\hookrightarrow I^{\bullet,b})\)。刚才的结论在每一步给出
\[
\begin{aligned}
0&\longrightarrow K_b^a\longrightarrow I^{a,b}
\longrightarrow K_{b+1}^a\longrightarrow0,\\
0&\longrightarrow B^a(K_b)\longrightarrow B_h^a(I^{\bullet,b})
\longrightarrow B^a(K_{b+1})\longrightarrow0,\\
0&\longrightarrow Z^a(K_b)\longrightarrow Z_h^a(I^{\bullet,b})
\longrightarrow Z^a(K_{b+1})\longrightarrow0,\\
0&\longrightarrow H^a(K_b)\longrightarrow H_h^a(I^{\bullet,b})
\longrightarrow H^a(K_{b+1})\longrightarrow0.
\end{aligned}
\]
例如最初两步是 \(0\to C\to I^{\bullet,0}\to K_1\to0\) 与 \(0\to K_1\to I^{\bullet,1}\to K_2\to0\)，纵微分 \(\partial_v^0\) 就是商到 \(K_1\) 后再包含进 \(I^{\bullet,1}\)。一般地，纵微分是复形态射
\(\partial_v^b:I^{\bullet,b}\twoheadrightarrow K_{b+1}\hookrightarrow I^{\bullet,b+1}\)，故与横微分交换。下一次纵微分先商掉 \(K_{b+1}\)，恰好杀掉这次的像，所以相邻复合严格为零。四类短正合列说明各增广纵列均正合，而其中的每项由引理内射，正是定义所需的四类内射消解。若 \(C\) 在 \(a<m\) 为零，上一引理的下界构造使所有 \(I^{\bullet,b},K_b\) 保留该下界。

任意 Cartan--Eilenberg 消解也具有刚才的逐行余核结构。在次数零，定义使增广在四类对象上都单；取余复形后，刚才的分析将其四类对象识别为相应余核。四类纵列的正合性又使这些余核单射到下一行，归纳得到相同的四类短正合列。因此以下比较适用于任意选定的消解。

给定 \(f:C\to D\)，这里用目标横行的内射性，沿源的单态射延拓；投射比较定理则用源的投射性，沿目标的满态射提升。先沿 \(C\hookrightarrow I_C^{\bullet,0}\)，把
\(C\to D\to I_D^{\bullet,0}\) 延拓为一个横向复形态射。它诱导第一余复形之间的态射；再沿 \(K_{C,1}\hookrightarrow I_C^{\bullet,1}\)，将
\(K_{C,1}\to K_{D,1}\to I_D^{\bullet,1}\) 延拓。每一步的嵌入在项、余边和上同调上都单，目标横行满足引理的内射性条件，故可以逐次延拓。余核的泛性质保证相邻延拓与纵微分交换，得到所需交换双复形映射。

为证明同伦唯一性，设两种提升为 \(f^b,g^b:I_C^{\bullet,b}\to I_D^{\bullet,b}\)，并记其差为 \(e^b\)。约定负消解次数的行及微分为零，取 \(s^0=0\)，归纳构造横向复形态射
\(s^b:I_C^{\bullet,b}\to I_D^{\bullet,b-1}\)。次数零的差杀掉增广 \(C\)，故经 \(K_{C,1}\) 因子化；沿 \(K_{C,1}\hookrightarrow I_C^{\bullet,1}\) 延拓后得到 \(s^1\)，使
\(e^0=s^1\partial_{v,C}^0\)。

若已在次数 \(b-1\) 得到
\[
e^{b-1}=\partial_{v,D}^{b-2}s^{b-1}
+s^b\partial_{v,C}^{b-1},
\]
令 \(h^b=e^b-\partial_{v,D}^{b-1}s^b\)。两个提升均与纵微分交换，故
\[
\begin{aligned}
h^b\partial_{v,C}^{b-1}
&=\partial_{v,D}^{b-1}e^{b-1}
-\partial_{v,D}^{b-1}s^b\partial_{v,C}^{b-1}\\
&=\partial_{v,D}^{b-1}\partial_{v,D}^{b-2}s^{b-1}=0.
\end{aligned}
\]
于是 \(h^b\) 杀掉 \(\operatorname{im}\partial_{v,C}^{b-1}=K_{C,b}\)，经余复形 \(K_{C,b+1}\) 因子化。沿
\(K_{C,b+1}\hookrightarrow I_C^{\bullet,b+1}\) 延拓为 \(s^{b+1}\)，便有
\[
e^b=\partial_{v,D}^{b-1}s^b+s^{b+1}\partial_{v,C}^b.
\]
这给出了全部纵同伦等式。

反交换化后，在双次数 \((a,b)\) 定义总同伦
\(S^{a,b}=(-1)^as^{a,b}\)。纵向两项给出
\(\partial_vs+s\partial_v\)，横向两项则为
\((-1)^a(\partial_hs-s\partial_h)=0\)，因为每个 \(s^b\) 都是横向复形态射。因此
\(f-g=\mathrm d_DS+S\mathrm d_C\)。总次数只有有限项，加性函子保持这些等式与有限双积，故作用任意加性函子后仍得到总链同伦。

对恒等复形取两个方向的比较映射，它们的复合与恒等总映射链同伦，因而给出与选择无关的总上同调。同伦 \(s^{a,b}\) 保持横次数 \(a\)，降低纵次数 \(b\)：对列滤过，它保持滤过，纵同伦使两种提升在第一页诱导同一映射；对行滤过，它降低滤过一层，不能直接保证第一页映射相同。不过每个 \(s^b\) 都与横微分交换，因而先取横上同调后，它诱导第一页纵复形的链同伦；再取纵上同调才得到第二页映射相同。以后各页由取同调确定。比较映射的复合性质同样由同伦唯一性保证，故第二页起的识别均自然。
\end{proof}

上面的显式嵌入是在复形层面使用内射性；它避免了逐个消解各项后还要修补 \(\mathrm{d}^2=0\) 的问题。消解本身通常很大，实际计算只会使用其横行分裂和四类纵列的消解性质。

\subsection{对象同调与总同调}
\label{b4:subsec:1502:03}

\begin{definition}\label{b4:def:hypercohomology}
设 \(F:\mathcal A\to\mathcal B\) 为交换范畴间加性左正合函子，\(\mathcal A\) 有足够内射对象，\(C\) 为 \(\mathcal A\) 中的有下界上链复形。取其 Cartan--Eilenberg 消解 \((I^{a,b},\partial_h,\partial_v)\)，以 \(\mathrm d_h=\partial_h\)、\(\mathrm d_v=(-1)^a\partial_v\) 反交换化。对 \(n\in\mathbb Z\)，定义
\[
 \mathbb H^n(F,C)=H^n\operatorname{Tot}(F(I)).
\]
这些对象称为 \(F\) 在 \(C\) 上的\term[chaoshangtongdiao]{超上同调}[hypercohomology]，也记作右超导出函子 \(\mathbb R^nF(C)\)。总化每次只有有限项，且上一比较定理保证选择独立及对复形映射的自然性。
\end{definition}
\symidx{\(\mathbb H^n(F,C),\ \mathbb R^nF(C)\)：函子在复形上的超上同调与右超导出函子}

\begin{theorem}\label{b4:thm:hypercohomology-spectral-sequences}
设 \(F:\mathcal A\to\mathcal B\) 为交换范畴间的加性左正合函子，\(\mathcal A\) 有足够内射对象，\(C\) 为 \(\mathcal A\) 中的有下界上链复形。则有两套自然的强收敛谱序列
\[
 {}^{\mathrm I}E_2^{p,q}=H^p(R^qF(C^\bullet))
 \Longrightarrow\mathbb H^{p+q}(F,C),
\]
\[
 {}^{\mathrm {II}}E_2^{p,q}=R^pF(H^q(C))
 \Longrightarrow\mathbb H^{p+q}(F,C).
\]
若 \(C\) 的每项都 \(F\)-零调，则 \(\mathbb H^n(F,C)\cong H^n(F(C))\) 对每个 \(n\in\mathbb Z\) 成立。对 \(\mathcal A\) 中的有下界上链复形 \(D\)，每个拟同构 \(C\to D\) 都诱导超上同调同构。
\end{theorem}
\begin{proof}
对 \(F(I)\) 使用\cref{b4:thm:double-complex-spectral-sequences}。纵列是 \(C^a\) 的内射消解，故纵向先取上同调给出 \(R^bF(C^a)\)。固定 \(b\) 后，这些对象以 \(R^bF(\mathrm d_C^a)\) 为微分组成横向复形：函子性使相邻复合为零，比较映射保证它正是由 \(\partial_h\) 诱导的微分。再取横上同调得到第一套第二页。纵微分只是通常消解微分乘常数 \((-1)^a\)，其核与像不变，故可直接用同一子商识别上同调。若改用逐项乘 \((-1)^{ab}\) 的复形同构作识别，则横微分相应多出 \((-1)^b\)；沿该横行再次作相同的符号调整便恢复通常微分，所得上同调相同。

横向先取上同调时，每行的余边、余圈和上同调均内射，两个基本短正合列分裂，所以加性 \(F\) 保持它们。微分分解为商到余边再包含进下一项，分裂性使这两张映射作用 \(F\) 后仍分别满、单，因此
\[
 \ker F(\partial_h^a)\cong F(Z_h^a(I^{\bullet,b})),\qquad
 \operatorname{im}F(\partial_h^{a-1})\cong F(B_h^a(I^{\bullet,b})).
\]
这给出典范同构
\[
 H_h^a(F(I^{\bullet,b}))\simeq F(H_h^a(I^{\bullet,b})).
\]
而 \(H^a(C)\to H_h^a(I^{\bullet,\bullet})\) 是内射消解，再取纵上同调便得到 \(R^bF(H^a(C))\)，交换两个指标即为第二式。其纵符号同样通过逐次乘 \((-1)^{ab}\) 识别，不改变结果。下界保证两种滤过逐度有限，故均强收敛。

若各 \(C^a\) 都零调，第一套谱序列只剩 \(q=0\) 一行。把横次数的下界平移至零后，单行退化的边缘同态给出 \(H^n(F(C))\cong\mathbb H^n(F,C)\)。若 \(C\to D\) 拟同构，第二套谱序列在第二页的各位置均为同构。两套消解的横次数各有下界，纵次数非负，因此每个固定总次数仅有有限多个双次数位置；相应滤过逐度有限且两边均强收敛。于是可从这一有限页应用\cref{b4:cor:spectral-comparison-finite}，得到目标的上同调同构。
\end{proof}

特别地，把对象 \(M\) 放在次数零时，\(\mathbb H^n(F,M[0])=R^nF(M)\)；若整个 \(C\) 正合，第二套第二页为零，超上同调全部消失。普通左正合函子却未必保持拟同构：取 \(F=\operatorname{Hom}_{\mathbb Z}(\mathbb Z/2\mathbb Z,-)\)，次数零、一上的 \(C=(\mathbb Z\hookrightarrow\mathbb Q)\) 到 \(D=(\mathbb Q/\mathbb Z)[-1]\) 的商映射是拟同构，但 \(F(C)=0\)，而 \(H^1(F(D))\cong\mathbb Z/2\mathbb Z\)。因此所做的进一步消解是在校正普通逐项函子破坏拟同构的问题，不能由一般左正合性推出上面的横向同调交换。内射上链版本的标准叙述见 Weibel\cite[Cohomology Variant 5.7.9]{Weibel1994HomologicalAlgebra}；相应投射链版本及其存在性、比较和同伦见同书\cite[Definition 5.7.1, Lemma 5.7.2, Exercises 5.7.2--5.7.3]{Weibel1994HomologicalAlgebra}。
